\section{Ameaças à Validade}
\label{sec:ameacas}

Nesta seção, são apresentadas as ameaças à validade deste estudo e as estratégias adotadas para mitigá-las, organizadas segundo a classificação de Wohlin et al.\ (2012)\nocite{wohlin:2012}.

Em relação à validade de construto, a ferramenta CK analisa o código-fonte estaticamente, sem resolver dependências externas compiladas, e pode subestimar heranças ou acoplamentos indiretos. Para mitigar essa ameaça, empregou-se a mesma versão fixa do extrator em todos os \emph{snapshots}, de modo que eventuais subestimações incidem igualmente sobre todas as eras, e valores indefinidos ou falhas nunca foram convertidos em zero. Além disso, o detector de IA captura apenas vestígios explícitos versionados de um catálogo fixo de agentes, que não cobre, por exemplo, correções automáticas de segurança nem \emph{bots} de revisão citados como coautores, e o grupo sem IA pode conter uso de assistentes que não deixou registro. Por isso, os nomes ``com IA'' e ``sem IA'' referem-se apenas à presença dos vestígios procurados, e a ausência de vestígio não é interpretada como ausência de uso. Por fim, os critérios de curadoria são heurísticas locais que podem classificar repositórios incorretamente; para reduzir esse risco, as regras foram definidas previamente, aplicadas de forma determinística e ajustadas às convenções do Spring Boot, com as decisões individuais registradas para auditoria.

Quanto à validade interna, as variações entre eras e os contrastes entre os grupos de IA coexistem com fatores concorrentes, como expansão do tamanho do código, refatorações e mudanças de equipe. Para mitigar essa ameaça, as comparações entre eras foram feitas sobre os mesmos repositórios, o tamanho do código (arquivos Java e LOC) foi reportado junto às métricas e nenhuma relação causal foi atribuída à IA. Os grupos de IA também diferem em atividade recente, pois 150 dos 449 projetos sem IA não receberam \emph{commits} na \emph{branch} principal desde 2022, ante nenhum dos 88 com IA, o que é reportado junto aos resultados. Ademais, a data do último envio de código é apenas uma aproximação de atividade, pois pode refletir \emph{tags} ou ramos secundários; por essa razão, a medição de E4 utiliza o estado efetivo da \emph{branch} principal na coleta, e os 150 repositórios sem \emph{commits} nessa \emph{branch} desde 2022 entram com variação nula em relação a E3, o que atenua as diferenças médias entre E3 e E4.

No que se refere à validade externa, os resultados podem não ser generalizáveis para todos os projetos Spring Boot, pois a seleção depende do tópico \texttt{spring-boot} atribuído no GitHub, da linguagem principal Java e da exigência de atividade a partir de 2022 e de medição em todas as eras, o que exclui projetos arquivados ou inativos e introduz viés de sobrevivência em favor de sistemas longevos. O grupo pioneiro, com 20 repositórios, é particularmente reduzido. Para mitigar essa ameaça, os dois grupos temporais foram analisados separadamente e as conclusões foram delimitadas ao conjunto observado.

Por fim, em relação à validade de conclusão, o estudo restringe-se à estatística descritiva, e as medições repetidas dos mesmos repositórios não constituem amostras independentes. Para mitigar essa ameaça, não foram reportados testes de hipótese nem tamanhos de efeito, cada repositório recebeu peso igual nas estatísticas, as médias foram acompanhadas de mínimo, mediana e máximo, e a exigência de medição em todas as eras eliminou a perda de repositórios entre elas, de modo que as diferenças reportadas correspondem exatamente às variações médias dos mesmos sistemas.
